\documentclass[aspectratio=169,11pt]{beamer}

\usetheme{Madrid}
\setbeamertemplate{navigation symbols}{}
\setbeamertemplate{footline}[frame number]

\usepackage{amsmath}
\usepackage{graphicx}
\usepackage{tikz}
\usepackage{xcolor}
\usetikzlibrary{arrows.meta,positioning,fit}

\definecolor{mainblue}{HTML}{1F5A8A}
\definecolor{accentgreen}{HTML}{2E8B57}
\definecolor{accentorange}{HTML}{D97706}
\definecolor{softgray}{HTML}{F4F6F8}
\definecolor{textgray}{HTML}{374151}

\setbeamercolor{frametitle}{bg=mainblue,fg=white}
\setbeamercolor{titlelike}{bg=mainblue,fg=white}
\setbeamercolor{structure}{fg=mainblue}
\setbeamercolor{block title}{bg=mainblue,fg=white}
\setbeamercolor{block body}{bg=softgray,fg=textgray}
\setbeamerfont{frametitle}{series=\bfseries}

\begin{document}

\begin{frame}[t]{Traffic Flow Generation}
  \footnotesize
  \begin{block}{Research Question}
    How can we learn a \textbf{controllable}, \textbf{realistic}, and \textbf{executable} traffic flow generator when traffic supervision is incomplete, e.g., sparse observations and missing labels?
  \end{block}

  \vspace{0.18cm}
  \centering
  \[
    \mathcal{G}_{\theta}:
    (\mathcal{R}, \mathcal{P}, \mathbf{c})
    \rightarrow
    \mathcal{F}_{\mathrm{sim}},
  \]
  \vspace{-0.1cm}
  \begin{columns}[T,onlytextwidth]
    \begin{column}{0.49\textwidth}
      \textbf{Input}
      \begin{itemize}
        \item $\mathcal{R}$: road network structure
        \item $\mathcal{P}$: POIs and urban context
        \item $\mathbf{c}$: time and scenario condition
      \end{itemize}
    \end{column}

    \begin{column}{0.49\textwidth}
      \textbf{Output}
      \begin{itemize}
        \item $\mathcal{F}_{\mathrm{sim}}$: executable traffic flow files for simulators
        \item Contains flows indexed by intersection $i$, road link $e$, movement $m$, and time $t$
      \end{itemize}
    \end{column}
  \end{columns}
\end{frame}

\begin{frame}[t]{Traffic Flow Generation}
  \footnotesize
  \begin{block}{Why It Matters}
    Signal control algorithms need realistic, controllable, and executable traffic demand
    for reliable evaluation, but real-world OD demand, routes, and turn movements are often
    missing or only sparsely observed.
  \end{block}

  \vspace{0.35cm}
  \textbf{Why labels are missing}
  \begin{itemize}
    \item Limited sensing coverage means many links and turns are never directly observed.
    \item Privacy and data-sharing constraints prevent some traffic data from being publicly released.
  \end{itemize}
\end{frame}

\begin{frame}[t]{Traffic Flow Generation}
  \footnotesize
  \begin{columns}[T,onlytextwidth]
    \begin{column}{0.48\textwidth}
      \textbf{Generation target}
      \begin{itemize}
        \item Realistic traffic flow files runnable by a simulator, consistent with real-world traffic patterns
        \item Controllable generation, allowing customization aligned with user preferences
      \end{itemize}
    \end{column}

    \begin{column}{0.49\textwidth}
      \textbf{Hierarchical generation:}
      \[
        \text{OD-time} \rightarrow \text{route}
        \rightarrow \text{link/turn} \rightarrow \text{flow}
      \]

      \begin{itemize}
        \item OD-time demand decomposition
        \item Route assignment
        \item Link/turn flow projection
        \item Simulation flow-file generation
      \end{itemize}
    \end{column}
  \end{columns}

  \vspace{0.3cm}
  \begin{block}{Main Contributions}
    \begin{enumerate}
      \item We study traffic-flow generation for signal control algorithms.
      \item We propose a hierarchical cascade from OD-time demand to routes, link/turn flows, and executable flow files.
      \item We validate the framework on real networks and traffic counts across Austin, Chicago, NYC, San Francisco, and Washington DC under zero-shot, few-shot retrieval, and full-shot settings.
    \end{enumerate}
  \end{block}
\end{frame}

\begin{frame}[t]{Traffic Flow Generation}
  \centering
  \vspace{-0.1cm}
  \IfFileExists{Figures/Overview.png}{%
    \includegraphics[width=\textwidth,height=0.84\textheight,keepaspectratio]{Figures/Overview.png}
  }{%
    \fbox{%
      \begin{minipage}[c][0.76\textheight][c]{0.92\textwidth}
        \centering
        Place the method figure here as \texttt{Figures/Overview.png}.
      \end{minipage}
    }
  }
\end{frame}

\begin{frame}[t]{Inference}
  \centering
  \vspace{-0.1cm}
  \IfFileExists{Figures/inference.png}{%
    \includegraphics[width=\textwidth,height=0.84\textheight,keepaspectratio]{Figures/inference.png}
  }{%
    \fbox{%
      \begin{minipage}[c][0.76\textheight][c]{0.92\textwidth}
        \centering
        Place the inference figure here as \texttt{Figures/inference.png}.
      \end{minipage}
    }
  }
\end{frame}

\end{document}
